\section{Introduction}\label{sec:intro}
Counterfactual analysis~\cite{wachter2017counterfactual,mothilal2020explaining} asks a central question in data-driven decision making: \textit{``what minimal and feasible change would lead to a different outcome?''} It seeks semantically meaningful modifications that alter a prediction and provide actionable alternatives. 
%By examining how outcomes shift under plausible changes, counterfactual reasoning supports explanation~\cite{dandl2020multi}, recourse~\cite{karimi2021algorithmic}, and auditing~\cite{lee2021auditing}, with important applications, such as healthcare~\cite{abrate2021counterfactual}, finance~\cite{yang2021exploring}, and law~\cite{zhang2023cfgl}. 
By examining how outcomes shift under plausible changes, counterfactual reasoning supports a wide range of tasks, including explanation~\cite{dandl2020multi}, recourse~\cite{karimi2021algorithmic}, and auditing~\cite{lee2021auditing}, and has been applied across diverse domains, including healthcare~\cite{abrate2021counterfactual}, finance~\cite{yang2021exploring}, and law~\cite{zhang2023cfgl}.
Recent work further leverages counterfactuals to improve diversity~\cite{mothilal2020explaining}, robustness~\cite{dominguez2022adversarial}, and fairness~\cite{von2022fairness}.

Despite this progress, most existing approaches adopt a fundamentally \emph{model-centric} perspective, framing the problem as \emph{how to modify the input to change a model's prediction}. Concretely, they analyze gradients, perturb features, or generate synthetic edits by interrogating a predictive model~\cite{bjoru2026framework,mothilal2020explaining,dandl2020multi,pawelczyk2020learning,kanamori2020dace,lucic2021focus,vanlooveren2020interpretable,galhotra2021explaining,schleich2021geco,kavouras2023fairness}. Such approaches are inherently tied to a specific model and often rely on decision boundaries or optimization objectives defined in representation space. As a result, the generated counterfactuals are frequently synthetic, may violate domain constraints, and can lack semantic validity or interpretability. The hypothetical changes constructed by such methods do not necessarily correspond to real data objects or feasible interventions~\cite{delaney2021uncertaintyestimationoutofdistributiondetection, mahajan2019preserving}.

We argue that, guided by feasibility, verifiability, and interpretability, counterfactuals in high-stakes settings~\cite{abrate2021counterfactual,yang2021exploring,lee2021auditing} should be grounded in observed data. This leads to a key question: \emph{which real alternatives explain the outcome?} Rather than synthesizing inputs, we identify existing instances with different outcomes that are \emph{close}, \emph{feasible}, and \emph{interpretable}. This \emph{data-centric view} is model-agnostic, supported by verifiable evidence, and casts counterfactual analysis as a first-class database query primitive over observed data.

We articulate a clear vision of \textbf{data-centric counterfactual analysis}: \emph{counterfactual reasoning should be elevated to a first-class data management primitive, where systems retrieve feasible, outcome-changing alternatives from observed data under explicit constraints, costs, and objectives}. Under this view, counterfactual reasoning is formulated as a constrained selection problem over observed data. The key novelty is to \textit{recast counterfactual (and causal) analysis as database query answering}, bringing it into the core of data management. Our vision extends the frontier of data management to meet the growing demands of AI by enabling structured, constraint-aware reasoning over data. Our vision is orthogonal to and complementary with recent causal graph-based approaches~\cite{pachera2025whatif}, which incorporate causal graphs to support interventions and counterfactual queries but rely on explicit causal models.

\rev{Our goal is not to develop another counterfactual generation algorithm, but to establish \emph{counterfactual querying} as a first-class data management primitive. The challenge is therefore not how to generate a counterfactual, but how to design the query abstraction, indexing mechanisms, optimization strategies, and execution model when counterfactual reasoning becomes a recurring database operation.

From this perspective, counterfactual querying is fundamentally different from existing retrieval pipelines. Existing systems compose vector search, filtering, reranking, and subset selection to retrieve similar objects. Counterfactual querying instead asks a different question: \emph{which feasible interventions would change the outcome?} Consequently, the opposite-outcome predicate and admissibility define the feasible search space, intervention cost $\omega$ replaces geometric similarity as the optimization objective, and the answer is a set of complementary counterfactuals that jointly optimizes a user-specified objective $\Omega$.

These semantics cannot be realized by simply composing existing retrieval operators. Treating retrieval, feasibility, intervention cost, and set selection as separate stages prevents holistic optimization and obscures the true semantics of the query. Our central thesis is therefore to make the opposite-outcome predicate, intervention cost $\omega$, and set objective $\Omega$ first-class components of the query model, enabling future data systems to optimize counterfactual reasoning as a native database operation.}



Realizing our vision requires addressing two key challenges: (i) how to support such queries efficiently and at scale, and (ii) how to enable principled reasoning across heterogeneous data modalities.

\subsection{Key Challenge 1: Efficient and Scalable Counterfactual Query Answering}

At first glance, counterfactual queries resemble similarity search, as both retrieve instances that are ``close'' to a query under some distance metric~\cite{chaves2001searching,indyk1998approximate,malkov2020efficient}. This suggests a simple view: counterfactuals can be obtained by finding nearest neighbors from a different class. Yet, this view is insufficient. In practice, users seek not merely proximity, but \emph{outcome-changing} alternatives that provide \emph{actionable insight}~\cite{lipton1990contrastive,miller2019explanation,shang2022whyami,bucinca2025contrastive}. This introduces a fundamentally contrastive requirement: retrieved instances must be close, have a different outcome, and remain feasible under domain constraints. Consequently, counterfactual querying goes beyond nearest-neighbor search—it must enforce \emph{label constraints}, respect \emph{admissibility conditions}, incorporate \emph{cost-faithful metrics aligned with real interventions}, and return \emph{informative rather than redundant answers}~\cite{cherumanal2025diversification}.

These requirements introduce fundamental challenges. Feasibility is governed by domain constraints that are often non-metric and cannot be captured by distance alone. Cost functions must reflect real-world intervention effort rather than geometric proximity, which breaks standard indexing assumptions. Moreover, selecting a small set of counterfactuals requires balancing objectives such as diversity and fairness, which cannot be achieved by independent nearest-neighbor search. Addressing these challenges calls for new designs in representation, indexing, constraint-aware retrieval, and set-level optimization.

\subsection{Key Challenge 2: Counterfactual Querying Across Modalities}
Real-world data is inherently multi-modal, combining tables, graphs, sequences, text, and images, each with distinct structures such as attributes, temporal dependencies, and graph topology. Existing systems are highly effective at learning shared representations across modalities, but this embedding-centric paradigm remains limited in its ability to \emph{query and reason} over such data in a structured way, such as composing information or enforcing constraints. This limitation becomes especially visible for counterfactual reasoning, where retrieval must respect structured feasibility conditions rather than similarity alone~\cite{wachter2017counterfactual}.

To address this, we advocate a data-centric approach that treats modalities as collections of entities, relationships, dependencies, and constraints. Under this view, counterfactual queries operate on \emph{semantic dimensions of change} rather than raw features. These dimensions vary across modalities, such as attributes in tables, temporal signals in sequences, or structures in graphs, and may appear in different forms across modalities. This heterogeneity makes similarity search alone insufficient, as it may return candidates that violate constraints or lack semantic validity. Set-level counterfactual optimization further adds complexity, since robustness~\cite{guyomard2023generating}, diversity~\cite{laugel2023achieving}, and fairness~\cite{kusner2017counterfactual} create dependencies among candidates. These challenges call for principled methods that jointly handle feasibility, interpretability, and set-level utility across modalities.

In summary, we present a data-centric vision of counterfactual analysis and formalize it as counterfactual querying, a first-class data management primitive with a query model, a two-phase framework, and core design principles. 

The remainder of the paper is organized as follows. Section~\ref{sec:query-answering} introduces the query model and framework. Section~\ref{sec:instance} studies instance-level semantics across data modalities, and Section~\ref{sec:set} presents set-level optimization for collective validity. Section~\ref{sec:systems} examines existing system limitations and presents a proof-of-concept system, DcCA. Section~\ref{sec:success} outlines success criteria, and Section~\ref{sec:con} concludes.

 


%\paragraph{Contributions.}
%This paper advocates counterfactual querying as a novel data management primitive and presents a vision for data-centric counterfactual analysis. We (i) formalize counterfactual queries as constrained, label-aware, and set-oriented retrieval over structured data, (ii) identify the fundamental gap between similarity search and counterfactual querying, highlighting the need for constraint-aware and cost-faithful retrieval, and (iii) develop a unified perspective for counterfactual reasoning across heterogeneous data modalities.

%The rest of the paper is organized as follows. Section~2 formalizes counterfactual queries and introduces the problem definition. Section~3 presents a two-phase framework for efficient and principled counterfactual querying. Section~4 analyzes instance-level semantics and system challenges across different data modalities. Section~5 studies set-level objectives and their role in shaping counterfactual results. Section~6 discusses evaluation criteria and open challenges, and Section~7 concludes the paper.

